\section{Test Case 18: high-resolution dynamic Maxwell evaporation}

Test Case 18 is the high-resolution counterpart of Test Case 16.  It uses the
Maxwell RK4 integrator with Yarin evaporation, nozzle insertion, and collector
removal enabled.  The initial jet length is 16 cm with \texttt{points 800},
giving a nominal spacing of 0.02 cm, that is 200 micrometres, between adjacent
nodes.  As in Test Case 16, dynamic refinement is disabled and the case is
serial.

Because the initial capacity already contains 800 intervals, the case measures
the dynamic topology path at a finer spatial resolution rather than forcing
frequent capacity reallocations.  As in Test Case 16, the OpenACC build
takes the common device step from the first step.

The NVFORTRAN CPU reference completed in about 14.2 s with 500 topology
additions, 899 removals, three reallocations, and 401 active beads.  Since
2026-10-07 an NVFORTRAN 24.3 run on an NVIDIA A30 gives the same totals,
every insertion at the CPU's step; before, its topology path inserted one
step later whenever the arrays had to grow, and the run gave 499 additions,
898 removals and 401 active beads.

The transverse components, of order $10^{-12}$~cm, still differ by up to
their own size, a sign of the case's sensitivity to the summation order of
the direct Coulomb force.  Test Case 18 is therefore a
performance and porting probe rather than a strict pointwise regression case,
and it is not part of the short smoke and regression matrices.  The input is
stored in \texttt{examples/input-18/input.dat}.
